\documentclass[11pt]{amsart}
\usepackage{geometry}                % See geometry.pdf to learn the layout options. There are lots.
\geometry{letterpaper}                   % ... or a4paper or a5paper or ... 
%\geometry{landscape}                % Activate for for rotated page geometry
%\usepackage[parfill]{parskip}    % Activate to begin paragraphs with an empty line rather than an indent
\usepackage{graphicx}
\usepackage{amssymb}
\usepackage{epstopdf}
\DeclareGraphicsRule{.tif}{png}{.png}{`convert #1 `dirname #1`/`basename #1 .tif`.png}

\title{PS 3.23}
%\author{The Author}
%\date{}                                           % Activate to display a given date or no date

\begin{document}
\maketitle
%\section{}
%\subsection{}
Resuelva el ejercicio anterior, pero considere que las cinco calificaciones del alumno se leerán dentro de un ciclo.\\

Datos: $MAT_1$, $CAL_{1.1}$, $CAL_{1,2}$, . . $CAL_{1,5}$, ... $MAT_{35}$, $CAL_{35.1}$,.  $CAL_{35.5}$ \\

Donde:\\

$MAT_i$  es una variable de tipo entero que representa la matrícula del alumno i ($1 \leq i \leq 35$).\\
$CAL{i,j}$ es una variable de tipo real que representa la calificación j del alumno i ($1 \leq j \leq 5,1 \leq i \leq 35$).

\end{document}  